\section{Preliminaries}
\label{sec:prelim}

We introduce notation and recall the building blocks used in our construction.
Throughout, $\mathbb{F}$ denotes a finite field of size $|\mathbb{F}| = p$ for a prime $p$, and $\lambda$ denotes the security parameter.
We write $\textsf{negl}(\lambda)$ for any function that is negligible in~$\lambda$, and $\ppt$ for probabilistic polynomial-time.

\subsection{Notation}
\label{sec:notation}

Matrices are denoted by bold uppercase letters $\mat{A}, \mat{B}, \mat{C} \in \mathbb{F}^{k \times k}$, and vectors by bold lowercase letters $\vect{x}, \vect{y}, \vect{z} \in \mathbb{F}^k$.
For a vector $\vect{x}$, we write $\vect{x}[i]$ for its $i$-th entry. For a
matrix $\mat{M}$, we write $\mat{M}[i,j]$ for the entry in the $i$-th row and $j$-th column. 
We denote the $i$-th row of $\mat{M}$ by $\mat{M}[i,*]$, and we will often write it simply as $\mat{M}[i]$. 
Similarly, we denote the $j$-th column of $\mat{M}$ by $\mat{M}[*,j]$. We write $[n] = \{1,2,\ldots,n\}$. For $t \in \mathbb{N}$, $I \sample [n]^t$ denotes the uniform sampling of $t$ indices from $[n]$. Similarly, $r \sample \mathbb{F}$ denotes uniform random sampling from
$\mathbb{F}$.



\subsection{Linear Error-Correcting Codes}
\label{sec:prelim-codes}

A \emph{linear code} over $\mathbb{F}$ is an $\mathbb{F}$-linear map
$\mathcal{C}: \mathbb{F}^k \to \mathbb{F}^n$ defined by a generator matrix
$\mat{G} \in \mathbb{F}^{k \times n}$, so that
$\mathcal{C}(\vect{m}) = \vect{m} \cdot \mat{G}$ for any input vector
$\vect{m} \in \mathbb{F}^k$.
The code has \emph{rate} $\rho = k/n$ and \emph{minimum distance} $d = \min_{\vect{c} \neq \vect{c}'} \Delta(\vect{c}, \vect{c}')$,
where the minimum is taken over distinct codewords
$\vect{c}, \vect{c}' \in \mathcal{C}$, and $\Delta$ denotes the Hamming distance.
The \emph{relative distance} is $\delta = d/n$.
A code is \emph{systematic} if the first $k$ coordinates of
$\mathcal{C}(\vect{m})$ equal $\vect{m}$ itself.

We write $\enc(\vect{m}) \coloneqq \mathcal{C}(\vect{m})$ for the encoding of an input vector.
For a matrix $\mat{M} \in \mathbb{F}^{k \times k}$, we define the
\emph{row-wise encoding} $\hat{\mat{M}} := \enc(\mat{M}) \in \mathbb{F}^{k \times n}$ as the matrix whose $i$-th row is $\enc(\mat{M}[i,*])$ for each $i \in [k]$.
The $j$-th column of $\hat{\mat{M}}$ is denoted by
$\hat{\mat{M}}[*,j] \in \mathbb{F}^k$ for each $j \in [n]$.

The key property we exploit is \emph{linearity}: for any coefficients
$r_1, \ldots, r_k \in \mathbb{F}$,
\begin{equation}\label{eq:code-linearity}
\sum_{i=1}^{k} r_i \cdot \hat{\mat{M}}[i,*]
=
\enc\!\left(\sum_{i=1}^{k} r_i \cdot \mat{M}[i,*]\right).
\end{equation}


\begin{definition}[Fold]\label{def:fold}
For vectors $\vect{c},\vect{r}\in \mathbb{F}^k$ with
$\vect{c}=(c_0,\ldots,c_{k-1})^\top$ and
$\vect{r}=(r_0,\ldots,r_{k-1})$, define
\[
\fold(\vect{c}, \vect{r}) \coloneqq \langle \vect{r}, \vect{c} \rangle = \sum_{i=0}^{k-1} r_i \cdot c_i.
\]
Equivalently, for the $j$-th column of a row-wise encoded matrix,
$\fold(\hat{\mat{M}}[*,j], \vect{r})$ equals the $j$-th coordinate of
$\enc(\vect{r}\mat{M})$.
This follows directly from the linearity of the code (Equation~\ref{eq:code-linearity}).
\end{definition}

\paragraph{Concrete instantiations}
Our protocol is compatible with any linear code.
Two natural choices are: (i) \emph{Reed--Solomon codes}, where $\mathcal{C}(\vect{m})$ evaluates the polynomial $\sum_i m_i X^i$ over a domain of size $n > k$, achieving the optimal (MDS) distance $\delta = 1 - k/n + 1/n$; and (ii) \emph{expander-based codes} as used in Brakedown~\cite{golovnev2023brakedown}, which achieve linear-time encoding $\mathcal{O}(n)$ with constant relative distance.
We discuss the trade-offs in Section~\ref{sec:complexity}.

% \subsection{Commitment Schemes}
% \label{sec:prelim-commitments}

% A commitment scheme lets a prover fix a value and later open it. We use
% commitments only through the standard binding property: after a commitment is
% fixed, a polynomial-time prover cannot open it to two different values except
% with negligible probability.

% A \emph{Merkle tree}~\cite{merkle1989certified} over a collision-resistant hash function $H: \{0,1\}^* \to \{0,1\}^\lambda$ provides a vector commitment scheme with the following algorithms:
% \begin{itemize}
%   \item $\cm \gets \merkleC(\vect{v})$: Given a vector $\vect{v} = (v_1, \ldots, v_n)$, construct a binary hash tree with leaves $H(v_1), \ldots, H(v_n)$ and output the root hash $\cm$.
%   \item $(\vect{v}_j, \pi_j) \gets \merkleO(\cm, j)$: Open position $j$, returning the leaf value $v_j$ and the authentication path $\pi_j$ (the sequence of sibling hashes from the leaf to the root).
%   \item $b \gets \merkleV(\pi_j, v_j, \cm)$: Verify that $v_j$ is the $j$-th leaf consistent with root $\cm$; output $b \in \{0,1\}$.
% \end{itemize}

% The scheme is \emph{position-binding}: under the collision resistance of $H$, no $\ppt$ adversary can produce $(j, v, \pi)$ and $(j, v', \pi')$ with $v \neq v'$ such that both verify against the same root $\cm$, except with probability $\epsilon_{\mathsf{bind}}(\lambda) = \textsf{negl}(\lambda)$.
\subsection{Commitment Schemes}
\label{sec:prelim-commitments}
A commitment scheme allows a committer to bind itself to a message while keeping the message hidden until it is opened. In this work, we use two types of commitments: Merkle tree commitments and Pedersen vector commitments.

\subsubsection{Merkle tree commitments}
A Merkle tree commitment over a collision-resistant hash function
$H: \{0,1\}^* \to \{0,1\}^\lambda$ is a vector commitment for a vector
$\vect{v}=(\vect{v}[1],\ldots,\vect{v}[n]) \in \mathbb{F}^n$.
It consists of the following algorithms.

\begin{itemize}
\item $\rt \gets \merkleC(\vect{v})$:
outputs a root hash $\rt$ obtained by constructing a binary hash tree with
leaves $H(\vect{v}[1]),\ldots,H(\vect{v}[n])$.

\item $(\vect{v}[j],\pi_j) \gets \merkleO(\vect{v},j)$:
outputs the leaf value $\vect{v}[j]$ and its authentication path $\pi_j$ for
position $j \in [n]$.

\item $b \gets \merkleV(\rt,j,\vect{v}[j],\pi_j)$:
outputs $b=1$ if $\vect{v}[j]$ is accepted as the $j$-th leaf consistent with
root $\rt$, and outputs $b=0$ otherwise.
\end{itemize}

Merkle tree commitments are \emph{position-binding}: under the collision resistance of $H$, no $\ppt$ adversary can produce a root $\rt$, an index $j \in [n]$, two distinct values $v \neq v'$, and authentication paths $\pi_j,\pi'_j$ such that
\[
\merkleV(\rt,j,v,\pi_j)= \merkleV(\rt,j,v',\pi'_j)=1,
\]
except with probability $\epsilon_{\mathsf{merkle}}(\lambda)=\mathsf{negl}(\lambda)$.



\subsubsection{Pedersen commitments}
Pedersen commitments can be viewed as algebraic vector commitments over a group $\mathbb{G}$ of prime order $q$.
$\mathsf{ck}=(G_1,\ldots,G_n,H)$ be a commitment key consisting of fixed generators in $\mathbb{G}$. A Pedersen commitment is defined by the following algorithms.

\begin{itemize}
\item $\cm \gets \pedC(\ck,\vect{v};r)$:
outputs a commitment $\cm := \sum_{i=1}^{n}\vect{v}[i]G_i + rH \in \mathbb{G}$ for a vector $\vect{v}\in\mathbb{F}^n$ and randomness $r\sample\mathbb{F}$.

\item $\pi \gets \pedO(\vect{v},r)$:
outputs the opening $\pi=(\vect{v},r)$ consisting of the committed vector and
the randomness.

\item $b \gets \pedV(\ck,\cm,\vect{v},r)$:
outputs $b=1$ if $\cm \stackrel{?}{=} \sum_{i=1}^{n}\vect{v}[i]G_i + rH$,
and outputs $b=0$ otherwise.
\end{itemize}

Pedersen commitments are \emph{perfectly hiding} and
\emph{computationally binding}: under the discrete logarithm assumption,
no $\ppt$ adversary can produce a commitment $\cm$, two distinct vectors
$\vect{v}\neq\vect{v}'$, and randomness values $r,r'$ such that
\[
\pedV(\ck,\cm,\vect{v},r) = \pedV(\ck,\cm,\vect{v}',r') = 1,
\]
except with probability
$\epsilon_{\mathsf{ped}}(\lambda)=\mathsf{negl}(\lambda)$.

They are also \emph{additively homomorphic}: for
$\cm_1=\pedC(\ck,\vect{v}_1;r_1)$ and
$\cm_2=\pedC(\ck,\vect{v}_2;r_2)$, $\cm_1+\cm_2 = \pedC(\ck,\vect{v}_1+\vect{v}_2;r_1+r_2).$

% \subsection{Succinct Non-interactive Argument of Knowledge}
% \label{sec:prelim-snark}

% \paragraph{SNARKs}
% A \emph{SNARK} (Succinct Non-interactive Argument of Knowledge) for an NP relation $\mathcal{R}$ consists of algorithms $(\textsf{Setup}, \snarkP, \snarkV)$:
% \begin{itemize}
%   \item $\para \gets \textsf{Setup}(1^\lambda, \mathcal{R})$: Generate public parameters.
%   \item $\pi \gets \snarkP(\para, \stmt, \wit)$: Given a statement $\stmt$ and witness $\wit$ with $(\stmt, \wit) \in \mathcal{R}$, produce a proof $\pi$.
%   \item $b \gets \snarkV(\para, \stmt, \pi)$: Verify the proof and output $b \in \{0,1\}$.
% \end{itemize}

% A SNARK satisfies \emph{completeness} (an honest prover always convinces the verifier), \emph{knowledge soundness} (any $\ppt$ prover that makes the verifier accept can be used to extract a valid witness), and \emph{succinctness} (the proof size $|\pi|$ and verifier time are $\textsf{poly}(\lambda, |\stmt|)$, independent of the witness size and computation).

% Our protocol is modular with respect to the choice of SNARK backend, including
% pairing-based systems such as Pinocchio~\cite{parno2013pinocchio} and
% Groth16~\cite{groth2016size}, as well as transparent or preprocessing systems
% such as Hyrax~\cite{wahby2018doubly}, Sonic~\cite{maller2019sonic},
% Marlin~\cite{chiesa2020marlin}, and Spartan~\cite{setty2020spartan}.
% The key metric is the \emph{circuit size} (number of R1CS constraints or PlonK gates) of the relation that the SNARK backend must prove, as this dominates the prover's cost.

% \paragraph{Commitment-carrying SNARKs}
% A \emph{commitment-carrying SNARK} (CC-SNARK) is a SNARK backend that lets the
% prover publish commitments $\Gamma_i$ to selected private witness values
% $a_i$ while proving an arithmetic relation that uses the same values.
% The public statement contains the commitments $\Gamma_i$.
% Soundness requires that, except with the SNARK backend's binding error, an
% accepting proof fixes each committed witness value: the value used by the
% relation at the selected witness slot is the value opened by $\Gamma_i$.
% In this paper, the commitments $\Gamma_{\chi,j}$ are SNARK-side commitments to
% sampled matrix columns or vector coordinates.

% \paragraph{Commit-and-prove SNARKs}
% A \emph{commit-and-prove SNARK} (CP-SNARK) proves an NP relation while also
% proving that designated witness values open public commitments.
% For a commitment scheme $\mathsf{Com}$ and relation $\mathcal{R}$, the induced
% commit-and-prove relation is
% \[
%   \begin{aligned}
%   \mathcal{R}_{\mathsf{CP}}
%   =
%   \{\,((\stmt,C),(v,\rho,\wit')) :\;&
%   C=\mathsf{Com}(v;\rho)\\
%   &{}\land(\stmt,v,\wit')\in\mathcal{R}\,\}.
%   \end{aligned}
% \]
% Thus a CP-SNARK binds an external commitment $C$ to the private value used in
% the proven computation.
\subsection{Succinct Non-interactive Arguments of Knowledge}
\label{sec:prelim-snark}

A SNARK (succinct non-interactive argument of knowledge) for an NP relation
$\relation$ is a tuple of algorithms $\bPi=(\setup,\prove,\verify)$.

\begin{itemize}
\item $\para \gets \bPi.\setup(1^\lambda,\relation)$:
generates public parameters $\para$ for the relation $\relation$.

\item $\pi \gets \bPi.\prove(\para,\stmt,\wit)$:
outputs a proof $\pi$ for a statement $\stmt$ and witness $\wit$ satisfying $(\stmt,\wit)\in\relation$.

\item $b \gets \bPi.\verify(\para,\stmt,\pi)$:
outputs $b=1$ if $\pi$ is accepted as a valid proof for $\stmt$, and outputs $b=0$ otherwise.
\end{itemize}

\subsection{CP-Link}
\label{sec:prelim-cplink}

A \emph{CP-Link} proof is a lightweight commit-and-prove proof whose only
purpose is to link two commitments to the same value.  The value may be a field
element, a vector, or a matrix column.
For a SNARK-side commitment $\Gamma$ and an external commitment $C$, the link
relation is
\[
  \begin{aligned}
  \mathcal{R}_{\mathsf{link}}
  =
  \{\,((\Gamma,C),(u,\rho_\Gamma,\rho_C)) :\;&
  \Gamma=\mathsf{Com}_{\mathsf{snark}}(u;\rho_\Gamma)\\
  &{}\land C=\mathsf{Com}_{\mathsf{ext}}(u;\rho_C)\,\}.
  \end{aligned}
\]
We use this relation in two places.  First, for each sampled position, it links
a SNARK-side commitment $\Gamma_{\chi,j}$ to an authenticated Merkle leaf handle
$L_{\chi,j}$.  Second, for each intermediate vector $v$, it links a SNARK-side
vector commitment $\Gamma^{\mathsf{vec}}_v$ to the pre-query vector commitment
$\mu_v$.  In our artifact, these equal-opening claims are batched into one
QA-NIZK proof.

\subsection{Freivalds' Algorithm}
\label{sec:prelim-freivalds}

Freivalds' algorithm~\cite{freivalds1977probabilistic} is a classical randomized algorithm for verifying matrix multiplication.
Given $\mat{A}, \mat{B}, \mat{C} \in \mathbb{F}^{k \times k}$, it samples $\vect{r} \sample \mathbb{F}^k$ and checks whether $\vect{r}(\mat{A}\mat{B}) = \vect{r}\mat{C}$.
The correctness relies on the following:

\begin{lemma}[Schwartz--Zippel~\cite{schwartz1980fast,zippel1979probabilistic}]\label{lem:sz}
Let $P(X_1, \ldots, X_m)$ be a non-zero polynomial of total degree $d$ over a finite field $\mathbb{F}$.
Then for $r_1, \ldots, r_m \sample \mathbb{F}$ chosen uniformly and independently,
\[
\Pr[P(r_1, \ldots, r_m) = 0] \le \frac{d}{|\mathbb{F}|}.
\]
\end{lemma}

In our protocol, we use a structured variant: instead of sampling
$\vect{r} \sample \mathbb{F}^k$ uniformly, we sample one scalar
$r \sample \mathbb{F}$ and set
$\vect{r} = (1, r, r^2, \ldots, r^{k-1})$.
This keeps the Freivalds challenge compact in the transcript and in the SNARK
statement: the verifier records one field element rather than $k$ independent
ones.
If $\mat{D} = \mat{A}\mat{B} - \mat{C} \neq \mat{0}$, then for any non-zero
column $\vect{d}_j$ of $\mat{D}$, the value
$\langle \vect{r}, \vect{d}_j \rangle$ is a non-zero univariate polynomial in
$r$ of degree at most $k-1$.
By Lemma~\ref{lem:sz}, it vanishes with probability at most
$(k-1)/|\mathbb{F}|$, which is negligible over the 256-bit fields used in our
experiments.
